\documentclass[11pt]{article}

\usepackage{fontspec}
\usepackage[ngerman]{babel}
\usepackage[a4paper,margin=2.6cm]{geometry}
\usepackage{amsmath,amssymb,amsthm}
\usepackage{microtype}
\usepackage{booktabs}
\setlength{\emergencystretch}{1.5em}

\newtheorem{satz}{Satz}[section]
\newtheorem{lemma}[satz]{Lemma}
\newtheorem{korollar}[satz]{Korollar}
\newtheorem{vermutung}[satz]{Vermutung}
\theoremstyle{definition}
\newtheorem{definition}[satz]{Definition}
\newtheorem{beispiel}[satz]{Beispiel}
\theoremstyle{remark}
\newtheorem{bemerkung}[satz]{Bemerkung}

\newcommand{\Hori}{\mathcal H}
\newcommand{\Struk}{\mathcal S}
\newcommand{\NN}{\mathbb N_0}
\newcommand{\val}{\operatorname{val}}
\newcommand{\ff}[2]{#1^{\underline{#2}}}
\newcommand{\KV}{K_V}
\newcommand{\KS}{K_S}

\title{Die Horimetrik:\\
Wert und Struktur in einem Zahlensystem\\
mit endlichen Fakultätshorizonten}
\author{Christian Felix Bürckert, M.\,Sc.\\[2pt]
\normalsize unter mathematischer Mitarbeit von Claude (Anthropic)\\
\normalsize und Ideenbeiträgen von GPT (OpenAI)}
\date{14. Juli 2026 --- Arbeitsstand, Version 0.16\\[4pt]
\normalsize Lizenz: CC\,BY\,4.0 --- Begleitprogramme: MIT\\
\normalsize\texttt{github.com/christianbuerckert/horimetrik}}

\begin{document}
\maketitle

\begin{abstract}
Wir untersuchen ein Mischbasen-Stellenwertsystem, dessen Basen zur
signifikantesten Stelle hin \emph{fallen}: Eine Ziffernfolge
$a_1\cdots a_n$ mit $0\le a_i\le i$ bezeichnet den Wert
$\sum_i a_i\,(n+1)!/(i+1)!$. Ein solches System ist nach links
versiegelt; jede Zahl lebt in einem endlichen \emph{Horizont} der
Größe $(n+1)!$, und eine führende Null verändert den Wert. Jedes
Element trägt dadurch zwei Identitäten: seinen \emph{Wert} und seine
\emph{Struktur} (ein Polynom mit nichtnegativen Koeffizienten in der
Basis der fallenden Fakultäten). Wir zeigen: (1)~Die wert- bzw.\
strukturtreuen Einbettungen der endlichen Horizonte erzeugen zwei
verschiedene direkte Grenzräume, $\NN$ und den positiven
Strukturhalbring. (2)~Der Raum trägt zwei spiegelbildliche
kommutative Halbringe; in beiden wirkt dieselbe ausgezeichnete Eins
als Projektor auf die jeweilige kanonische Kopie. (3)~Ein
\emph{Seitentreue-Satz}: Keine kommutative Halbring-Arithmetik mit
wert- bzw.\ strukturgetragenen Operationen kann die beiden
Identitäten mischen, unabhängig von der Wahl der Horizontregeln;
das Hindernis ist genau die Ziffernschranke des Systems. Ergänzend
beweisen wir ein asymptotisches Gesetz für den Kommutator der beiden
Kompaktifizierungen, ein Extremalprinzip mit exakten
Wachstumsgesetzen des Strukturhorizonts unter Addition,
Multiplikation, endlicher Differenz und Komposition (der
Kompositionsabschluss des positiven Halbrings selbst ist Teil des
Resultats), und dokumentieren neun offenbar unbekannte Zählfolgen. Der Text erhebt keinen Neuheitsanspruch über einen
dokumentierten Suchbefund hinaus.
\end{abstract}

\section{Die Idee}

Gewöhnliche Stellenwertsysteme lassen die signifikanteste Stelle am
weitesten zählen. Dieser Text untersucht die umgekehrte Konvention.
Für festes $n$ bezeichne eine Ziffernfolge $a_1a_2\cdots a_n$ mit
$0\le a_i\le i$ den Wert
\[
V_n(a)=\sum_{i=1}^n a_i\,\frac{(n+1)!}{(i+1)!}\,;
\]
die linke Stelle ist binär, trägt aber das größte Gewicht. Zwei
Konsequenzen: Das System ist \emph{nach links versiegelt} --- unter
der Basis $2$ gibt es nichts, mit $n$ Stellen sind genau $(n+1)!$
Werte darstellbar, und mehr Platz erfordert ein neues
Koordinatensystem. Und \emph{führende Nullen sind semantisch}:
$12\mapsto012\mapsto0012$ bezeichnet nacheinander $5,6,7$.

Jede Zahl trägt hier also zweierlei: einen Wert und eine Gestalt.
Die Mathematik dieses Doppelcharakters ist der Gegenstand des
Textes. Die Idee entstand aus der Frage, was passiert, wenn man die
Exponentenordnung einer Zahldarstellung umdreht; die Endlichkeit der
Horizonte ist dabei kein Designziel, sondern eine Zwangsläufigkeit.

\section{Der Horiraum}

\begin{definition}
Für $n\ge0$ sei $H_n=\{(a_1,\dots,a_n):0\le a_i\le i\}$, also
$|H_n|=(n+1)!$, mit der Wertabbildung $V_n$ wie oben und der
Codierung $E_n=V_n^{-1}$. Ferner sei
$\mu(x)=\min\{n:x<(n+1)!\}$ der \emph{Minimalhorizont} von
$x\in\NN$.
\end{definition}

\begin{satz}\label{satz:bij}
$V_n\colon H_n\to\{0,\dots,(n+1)!-1\}$ ist bijektiv.
\end{satz}

\begin{proof}
Induktion über $n$ mittels
$V_n(a)=(n+1)V_{n-1}(a_1,\dots,a_{n-1})+a_n$ und $0\le a_n\le n$:
Division mit Rest durch $n+1$ liefert Eindeutigkeit; der Wertebereich
ist $[0,(n+1)!-1]$, und beide Seiten haben Mächtigkeit $(n+1)!$.
\end{proof}

\begin{definition}
Mit $\ff Xd=X(X-1)\cdots(X-d+1)$ sei
$\Struk=\bigl\{\sum_d c_d\ff Xd:c_d\in\NN\bigr\}$ und
$\sigma\bigl(\sum c_d\ff Xd\bigr)=\max_{c_d>0}(d+c_d)$,
$\sigma(0)=0$. Das \emph{Strukturpolynom} von $a\in H_n$ ist
$P_a=\sum_{d=0}^{n-1}a_{n-d}\ff Xd$; es gilt $V_n(a)=P_a(n+1)$, und
$P\in\Struk$ ist genau dann in $H_n$ darstellbar, wenn
$\sigma(P)\le n$. Der \emph{Horiraum} ist
$\Hori=\{(n,P):P\in\Struk,\ n\ge\sigma(P)\}$, gleichwertig
$\{(n,x):n\ge\mu(x)\}$.
\end{definition}

$\Struk$ ist unter Addition und Multiplikation abgeschlossen (die
Produktformel der fallenden Fakultäten hat nichtnegative
Strukturkonstanten) und damit ein kommutativer Halbring.

\section{Erstes Hauptresultat: zwei Grenzräume}

Es gibt zwei natürliche Arten, einem Element mehr Platz zu geben:
die \emph{strukturtreue Erweiterung} $Z(n,P)=(n+1,P)$ (führende
Null; die Gestalt bleibt, der Wert wird zu $P(n+2)$) und der
\emph{werttreue Lift} $L(n,P)=(n+1,E_{n+1}(P(n+1)))$ (der Wert
bleibt, die Gestalt wird neu codiert). Im Dezimalsystem fallen beide
zusammen; hier nicht.

\begin{satz}\label{satz:limits}
Als gerichtete Systeme gilt
\[
\varinjlim(H_n,Z_n)\cong\Struk,
\qquad
\varinjlim(H_n,L_n)\cong\NN .
\]
\end{satz}

\begin{proof}
Unter $Z$ werden genau Ziffernfolgen identifiziert, die sich um
führende Nullen unterscheiden; jede Klasse trägt genau ein
Strukturpolynom, und jedes $P\in\Struk$ tritt (ab $H_{\sigma(P)}$)
auf. Unter $L$ gilt $L_{m,k}\circ L_{n,m}=L_{n,k}$ wegen
$V\circ E=\mathrm{id}$, und zwei Elemente werden genau dann
identifiziert, wenn ihre Werte übereinstimmen; jeder Wert tritt auf.
\end{proof}

Dieselben endlichen Räume erzeugen also, je nach Konsistenzbegriff
zwischen den Horizonten, zwei verschiedene unendliche Welten:
die Werte $\NN$ und die Strukturen $\Struk$. Alles Weitere handelt
vom Verhältnis dieser beiden Welten.

\begin{bemerkung}[dritte Erweiterung; Cantor-Koordinate]
Es gibt eine dritte natürliche Erweiterung: das \emph{Anhängen}
einer Null rechts, $R(a_1,\ldots,a_n)=(a_1,\ldots,a_n,0)$. Wegen
$V_n(a)=(n+1)!\cdot F(a)$ mit dem Fakultätsbruch
$F(a)=\sum_i a_i/(i+1)!\in[0,1)$ ist $R$ \emph{bruchtreu} und
skaliert den Wert exakt mit $n+2$; der zugehörige Grenzraum ist die
Menge der endlichen Cantor-Fakultätsbrüche. Der naheliegende Einwand
,,im Kern seit Cantor bekannt`` trifft damit genau diese dritte,
bruchtreue Welt -- nicht die gestalttreue (deren Grenzraum der
Strukturhalbring ist) und nicht die im Folgenden bewiesenen Sätze
über das Zusammenspiel. Bruch, Gestalt und Wert sind paarweise
verschiedene Invarianten.
\end{bemerkung}

\section{Zweites Hauptresultat: duale Projektorhalbringe}

Die beiden \emph{Kompaktifizierungen}\footnote{Nicht
topologisch gemeint: idempotente Retraktionen auf den jeweiligen
Minimalhorizont.}
$\KV(n,x)=(\mu(x),x)$ (Wert behalten, kleinstes Zuhause) und
$\KS(n,P)=(\sigma(P),P)$ (Gestalt behalten, kleinstes Zuhause) sind
idempotent und kommutieren nicht; das kleinste Beispiel ist die
Ziffernfolge $010$ in $H_3$ (Wert 4).

\begin{satz}[Wertseite]\label{satz:wertseite}
Mit
\[
(n,x)\oplus(m,y)=\bigl(\max(n,m,\mu(x{+}y)),\,x{+}y\bigr),
\qquad
(n,x)\otimes(m,y)=\bigl(\mu(xy),\,xy\bigr)
\]
ist $\Hori$ ein kommutativer Halbring ohne Eins; $(0,0)$ ist neutral
und absorbierend; die $\KV$-Fixpunkte bilden einen zu
$(\NN,+,\cdot)$ isomorphen Unterhalbring mit Eins $(1,1)$, und für
alle $h$ gilt
\[
h\otimes(1,1)=\KV(h).
\]
\end{satz}

\begin{proof}
Kern ist die Monotonie von $\mu$. Assoziativität von $\oplus$: Beide
Klammerungen liefern
$\max(n,m,k,\mu(\text{Teilsumme}),\mu(x{+}y{+}z))$, und die inneren
$\mu$-Terme werden von $\mu(x{+}y{+}z)$ absorbiert. Distributivität:
$\max(\mu(xy),\mu(xz),\mu(xy{+}xz))=\mu(x(y{+}z))$, wieder Monotonie.
Neutralität von $(0,0)$ folgt aus $n\ge\mu(x)$. Ein
$\otimes$-neutrales Element müsste $\mu(x)=n$ für alle Elemente
erzwingen, was an $(1,0)$ scheitert. Schließlich
$(n,x)\otimes(1,1)=(\mu(x),x)=\KV(n,x)$.
\end{proof}

\begin{lemma}\label{lem:sigmamono}
$\sigma(P+Q)\ge\max(\sigma(P),\sigma(Q))$.
\end{lemma}

\begin{proof}
An der Maximalstelle $d^*$ von $P$ gilt
$c_{d^*}(P{+}Q)\ge c_{d^*}(P)>0$.
\end{proof}

\begin{satz}[Strukturseite, Dualität]\label{satz:dualitaet}
Mit
\[
(n,P)\boxplus(m,Q)=\bigl(\max(n,m,\sigma(P{+}Q)),\,P{+}Q\bigr),
\qquad
(n,P)\boxtimes(m,Q)=\bigl(\sigma(PQ),\,PQ\bigr)
\]
ist $\Hori$ ebenfalls ein kommutativer Halbring ohne Eins; die
$\KS$-Fixpunkte bilden einen zu $\Struk$ isomorphen Unterhalbring;
$\KS$ ist Homomorphismus beider Operationen, und für alle $h$ gilt
\[
h\boxtimes(1,1)=\KS(h).
\]
\end{satz}

\begin{proof}
Wörtlich spiegelbildlich zu Satz~\ref{satz:wertseite}, mit
Lemma~\ref{lem:sigmamono} anstelle der $\mu$-Monotonie. Die
Homomorphie von $\KS$: für $\boxplus$ ist
$\max(\sigma P,\sigma Q,\sigma(P{+}Q))=\sigma(P{+}Q)$ nach dem
Lemma; für $\boxtimes$ sind beide Seiten
$(\sigma(PQ),PQ)$.
\end{proof}

\begin{bemerkung}
Es ist \emph{dieselbe} Hori-Zahl $(1,1)$ --- die Eins in ihrem
kleinsten Zuhause ---, die auf der Wertseite $\KV$ und auf der
Strukturseite $\KS$ als innere Multiplikation trägt. Das
Nichtkommutieren der Kompaktifizierungen ist das Nichtkommutieren
zweier Multiplikationen mit ein und demselben Element. Das von
$\KS,\KV$ erzeugte Transformationsmonoid ist übrigens endlich: Es
hat genau sechs Elemente mit der Relation
$\KV\KS\KV=\KS\KV$ (via eines Stabilitätslemmas: $\KS$ erhält
Wertkompaktheit) und fünf Idempotenten.
\end{bemerkung}

\section{Drittes Hauptresultat: der Seitentreue-Satz}

Eine Operation auf $\Hori$ heiße \emph{wertgetragen}, wenn der
Ergebniswert Summe bzw.\ Produkt der Operandenwerte ist (der
Ergebnishorizont ist frei, sofern gültig), und
\emph{strukturgetragen}, wenn das Ergebnispolynom Summe bzw.\
Produkt der Operandenpolynome ist (Horizont frei $\ge\sigma$).
Sätze~\ref{satz:wertseite} und~\ref{satz:dualitaet} zeigen: Beide
\emph{seitenreinen} Kombinationen sind mit geeigneten Horizontregeln
gesetzestreu. Die gemischten sind es nie:

\begin{satz}[Seitentreue]\label{satz:seitentreue}
\emph{(A)} Ist $\oplus$ wertgetragen und $\otimes$ strukturgetragen,
so ist bereits die Links-Distributivität für jede Wahl der
Horizontregeln unerfüllbar.
\emph{(B)} Ist $\oplus$ strukturgetragen und $\otimes$ wertgetragen,
so sind Links-Distributivität und Kommutativität von $\otimes$
zusammen unerfüllbar.
\end{satz}

\begin{proof}
\emph{(A).} Sei $u=(1,1)$ (Struktur $1$) und $e_1=u$,
$e_m=e_{m-1}\oplus u$, also $\val(e_m)=m$; die Struktur $Q_m$ von
$e_m$ ist irgendeine gültige Darstellung von $m$.
Links-Distributivität liefert induktiv
\begin{equation}\label{eq:kette}
\val(a\otimes e_m)=m\cdot\val(a\otimes u).
\end{equation}

\emph{Schritt 1: jedes $Q_m$ ist konstant.} Für $c\ge1$ hat
$a_c=(c,c)$ die konstante Struktur $c$; Konstanten werten überall zu
sich selbst aus, also $\val(a_c\otimes u)=c$. Mit \eqref{eq:kette}
folgt $c\,Q_m(y_c{+}1)=mc$, also $Q_m(y_c{+}1)=m$, wobei
$y_c\ge\sigma(cQ_m)=\max_d(d+c\,q_d)\ge c$ unbeschränkt wächst. Alle
Auswertungspunkte liegen im streng monotonen Bereich von $Q_m$
(Argumente $\ge\sigma+1\ge\deg+2$); bei $\deg Q_m\ge1$ gäbe es
höchstens einen solchen Punkt. Also ist $Q_m$ konstant, und mit
$a=u$ in \eqref{eq:kette} folgt $Q_m=m$.

\emph{Schritt 2: Konstanten sterben.} Sei $g=(2,3)$ mit Struktur
$X$ und $r=\val(g\otimes u)=t_g{+}1\ge3$ fest. Nach Schritt~1 hat
$g\otimes e_m$ die Struktur $mX$ mit $\sigma(mX)=m{+}1$, also
$\val(g\otimes e_m)=m(T{+}1)\ge m(m{+}2)$; nach \eqref{eq:kette}
ist dieser Wert $mr$ --- Widerspruch für $m>r-2$.

\emph{(B).} Nun addieren sich Strukturen und multiplizieren sich
Werte. Sei $f_1=u$, $f_m=f_{m-1}\oplus u$; dann hat $f_m$ die
konstante Struktur $m$ und den Wert $m$. Sei $A$ die Struktur von
$g\otimes u$ (eine Darstellung der $3$). Induktiv hat $g\otimes f_m$
die Struktur $mA$; als $\otimes$-Ergebnis vom Wert $3m$, realisiert
auf seinem Horizont $T_m$, erzwingt die Gültigkeit (Ziffernschranke)
$T_m\ge\sigma(mA)\ge m$ und $A(T_m{+}1)=3$. Wie in Schritt~1 ist
$A$ konstant, also $A=3$. Betrachte nun $f_2'=g\oplus g$ mit
Struktur $2X$, Horizont $H\ge\sigma(2X)=3$ und Wert
$2(H{+}1)\ge8$. Links-Distributivität und Kommutativität geben
\[
u\otimes(g\oplus g)=(g\otimes u)\oplus(g\otimes u),
\]
rechts mit Struktur $2\cdot3=6$, also Wert $6$; links steht der Wert
$2(H{+}1)\ge8$. Widerspruch.
\end{proof}

\begin{bemerkung}
Beide Widersprüche benutzen dieselbe Eigenschaft:
$\sigma(m\cdot P)$ wächst linear in $m$, weil die Ziffernschranke
$c_d\le n-d$ große Koeffizienten in große Horizonte zwingt. Die
Starrheit der beiden Identitäten ist damit eine direkte Konsequenz
der definierenden Ziffernbedingung $0\le a_i\le i$. Auch der
nichtkommutative Mischfall vom Typ~(B) ist unter beidseitiger
Distributivität ausgeschlossen: Die Rechts-Distributivität spiegelt
die Kettenkonstruktion und ersetzt die Kommutativität. Und die
Seitentreue ist inzwischen \emph{dreiseitig}: Die Bruch-Identität
(Cantor-Koordinate) trägt keine eigene Arithmetik -- ihre Addition
kann die Eins überschreiten und ist nie total -- und lässt sich mit
keiner der beiden Seiten verweben (Wert-$\oplus$ scheitert am
konstanten Fakultätsquotienten, Struktur-$\oplus$ an der
Halbkoeffizienten-Obstruktion $2A(Y)=Y^{\underline k}$).
Es gibt genau zwei Rechenwelten.
\end{bemerkung}

\section{Das Tiefengesetz des Kommutators}

Der Defekt
$\delta(h)=\operatorname{hor}(\KS\KV h)-\operatorname{hor}(\KV\KS h)$
ist beidseitig unbeschränkt, aber extrem asymmetrisch: Nach oben
tritt $+2$ erstmals bei $h=(9,720)$ auf ($720=6!=10\cdot9\cdot8$ hat
in $H_9$ die Struktur $\ff X3$; exhaustiv gesichert); die weiteren
Schwellen sind nur nach oben beschränkt: $+3$ spätestens ab $n=14$,
$+4$ spätestens ab $n=20$, $+5$ spätestens ab $n=26$.
Nach unten ist die Tiefe $-g$ an die \emph{Schalenlücke} gebunden;
die Sprungstellen sind
\[
m_g=3,\ 15,\ 84,\ 533,\ 3883,\ 31998,\ 294875\qquad(g=1,\dots,7).
\]

\begin{satz}\label{satz:asymptotik}
Mit $c_g=(g{+}2)!\sum_{k\ge1}k/(g{+}k{+}1)!=1+O(1/g)$ gilt
\[
\frac{(g+2)!}{c_g}\le m_g+1\le(g+2)!,
\qquad\text{insbesondere}\quad
\frac{m_g}{(g+2)!}\to1 .
\]
\end{satz}

\begin{proof}
Umindexieren liefert exakt
$M_{m-g}(m{+}1)/m!=(m{+}1)\sum_{k=1}^{m-g}k/(g{+}k{+}1)!$, wobei
$M_s=\sum_{d<s}(s{-}d)\ff Xd$ das Maximalpolynom ist. Die rechte
Seite wächst streng in $m$; für $m{+}1\ge(g{+}2)!$ genügt bereits
der Term $k{=}1$, und für $m{+}1<(g{+}2)!/c_g$ ist auch die volle
Reihe zu klein. Die Schranke $c_g\le1+\sum_{k\ge2}k(g{+}3)^{1-k}$
gibt $c_g=1+O(1/g)$.
\end{proof}

Inzwischen ist auch die Gleichheit bewiesen: Das Minimum von
$\delta$ über die Schale $m$ ist \emph{exakt} $-g(m)$, mit der
Terminaldarstellung von $M_{s_{\min}}(m{+}1)$ als stets tauglichem
(im Allgemeinen nicht einzigem) Zeugen.
Empirisch ist $(g{+}2)!/c_g$ sogar auf $\pm1$ genau gleich $m_g$;
die daraus abgeleitete Vorhersage $m_7=294\,875$ wurde
\emph{exakt} bestätigt, ebenso zwei frühere Vorhersagen dieses
Gesetzes ($-3$ ab Schale ${\approx}100$; $m_6=31998$ bei
prognostizierten ${\approx}31930$).

\section{Das Horizontkalkül}

Für einen $k$-stelligen Operator $T$ auf dem Strukturhalbring
$\Struk$ sei
\[
\Gamma_T(r_1,\ldots,r_k)=\max\{\sigma(T(P_1,\ldots,P_k)):
\sigma(P_i)\le r_i\}
\]
der maximale Strukturhorizont des Ergebnisses.
Wir schreiben $P\preceq Q$ für die koeffizientenweise Ordnung in der
fallenden Fakultätsbasis und
$M_r(X)=\sum_{d<r}(r{-}d)\,\ff Xd$ für das maximal gefüllte Polynom
mit $\sigma(M_r)=r$.

\begin{satz}[Extremalprinzip]\label{satz:extremal}
Ist $T$ in jedem Argument koeffizientenmonoton --- was auf jeden aus
Addition, Multiplikation und $\Delta$ zusammengesetzten Operator
zutrifft, da alle Strukturkonstanten nichtnegativ sind ---, so gilt
\[
\Gamma_T(r_1,\ldots,r_k)=\sigma\bigl(T(M_{r_1},\ldots,M_{r_k})\bigr).
\]
\end{satz}

\begin{proof}
$\sigma(P)\le r$ ist äquivalent zu $P\preceq M_r$ (Ziffernschranke
$c_d\le r-d$), und $\sigma$ ist $\preceq$-monoton; Monotonie von $T$
liefert die obere Schranke, $M_{r_i}\in\mathcal F_{r_i}$ die
Annahme.
\end{proof}

Damit besitzt jede Grundoperation ihr eigenes exaktes
Wachstumsgesetz: $\Gamma_+(r,s)=r+s$ (linear, scharf bei
Konstanten), $\Gamma_\times(r,s)=\beta(r,s)=\sigma(M_rM_s)$ (im
Wesentlichen fakultätisch, siehe unten) --- und für die endliche
Differenz:

\begin{satz}[Differenzüberlauf]\label{satz:differenzueberlauf}
Für $r\ge2$ gilt
$\Gamma_\Delta(r)=\lfloor(r+1)^2/4\rfloor-1$, angenommen bei der
Ein-Ziffern-Struktur $P=(r{-}d)\,\ff Xd$, $d=\lfloor(r+1)/2\rfloor$.
\end{satz}

\begin{proof}
Nach Satz~\ref{satz:extremal} ist
$\Gamma_\Delta(r)=\sigma(\Delta M_r)$ mit Koeffizienten
$b_k=(k{+}1)(r{-}k{-}1)$, also
$\Gamma_\Delta(r)=\max_{0\le k\le r-2}(k{+}1)(r{-}k)-1$; das Produkt
zweier Zahlen mit Summe $r{+}1$ ist maximal
$\lfloor(r{+}1)^2/4\rfloor$. (Exhaustiv bestätigt bis $r=8$:
$1,3,5,8,11,15,19$.)
\end{proof}

$\Delta$ senkt also den Grad und erhöht die Kapazität quadratisch:
Der Grad misst Dimension, $\sigma$ misst Kapazität; eine
Gradfiltration sieht diesen Effekt nicht.

Das Kalkül erfasst auch das Einsetzen:

\begin{satz}[Kompositionsabschluss]\label{satz:komposition}
$\Struk$ ist unter Komposition abgeschlossen, und $\circ$ ist in
beiden Argumenten koeffizientenmonoton; insbesondere
$\Gamma_\circ(r,s)=\sigma(M_r\circ M_s)$.
\end{satz}

\begin{proof}[Beweisskizze]
Wegen $P\circ Q=\sum_d c_d(P)(Q)^{\underline d}$ genügt
$(Q)^{\underline d}\in\Struk$. Interpretiert man $Q(x)$ als Anzahl
der Paare aus einer Schablone der Stelligkeit $k$ (je $a_k$ Stück)
und einer Injektion $[k]\hookrightarrow[x]$, so zählt
$Q(x)^{\underline d}$ die $d$-Tupel paarweise verschiedener solcher
Paare. Klassifikation nach der Bildvereinigung liefert
$Q(x)^{\underline d}=\sum_m N_m\binom xm$; auf den ausschöpfenden
Tupeln wirkt $S_m$ durch Umbenennung der Grundmenge \emph{frei}
(eine Permutation, die jede Injektion fixiert, fixiert die
ausgeschöpfte Vereinigung punktweise), also $m!\mid N_m$ und
$(Q)^{\underline d}=\sum_m(N_m/m!)X^{\underline m}\in\Struk$ ---
Positivität und Ganzzahligkeit in einem Zug. Monotonie in $Q$ folgt,
weil größere Schablonenvorräte die Tupelmengen nur vergrößern.
(Exhaustiv bestätigt für $r,s\le4$.)
\end{proof}

Die Diagonale $\Gamma_\circ(r,r)=1,4,23,6541,477149751,\ldots$
wächst turmartig und ist die neunte OEIS-unbekannte Folge des
Projekts; ihre Asymptotik ist offen. Der Verschiebeoperator
$E\,P(X)=P(X{+}1)$ als einfachster Spezialfall kostet exakt
$\Gamma_E(r)=r+\lfloor r^2/4\rfloor=\Gamma_\Delta(r{+}1)$ ---
Verschieben kostet so viel wie Ableiten einen Horizont höher; diese
Wertfolge ist die klassische Viertelquadrat-Folge (A024206), neu ist
das Gesetz, nicht die Folge. Das stetige Analogon des
Kompositionsabschlusses ist klassisch (absolut monotone Funktionen,
Faà di Bruno); Satz~\ref{satz:komposition} ist die diskrete
Gitterverschärfung mit der zusätzlichen Teilbarkeit $m!\mid N_m$. Additiv ist
$(\mathcal F_r)_r$ mit $\mathcal F_r=\{P:\sigma(P)\le r\}$,
$|\mathcal F_r|=(r{+}1)!$, eine klassisch verträgliche Filtration;
multiplikativ ist sie nichtlinear mit optimalem Wachstumsgesetz
$\beta$. Via Newton-Entwicklung ($c_d=\Delta^dP(0)/d!$) ist
$\sigma(P)=\max_{\Delta^dP(0)>0}\bigl(d+\Delta^dP(0)/d!\bigr)$ eine
basisfreie \emph{Kapazitätshöhe}; sie ist keine Bitkomplexität
($\sigma$ einer Konstanten $c$ ist $c$ selbst).

\section{Die Zählbreite: eine \#P-Interpretation}

Wegen $X^{\underline d}=d!\binom Xd$ hat jedes $P\in\Struk$
nichtnegative ganzzahlige Koeffizienten in der Binomialbasis
($b_d=c_d\,d!$) und liegt damit in der Klasse der
\emph{binomial-good}-Polynome, die als die (relativierenden)
polynomiellen \#P-Abschlussoperationen charakterisiert sind
(Ikenmeyer--Pak, FOCS 2022, arXiv:2204.13149). $\Struk$ ist eine
\emph{echte} Teilklasse ($\binom X2\notin\Struk$); die
$d!$-Teilbarkeit entspricht geordneten Tupeln statt ungeordneten
Auswahlen.

\begin{satz}[Hori-\#P-Satz]\label{satz:sharpp}
Für $f\in\#\mathrm P$ und $P\in\Struk$ ist
$P\circ f\in\#\mathrm P$.
\end{satz}

\begin{proof}
Rate $d\le\deg P$, eine Variante $\ell\le c_d$ und $d$
Zeugenkandidaten; akzeptiere bei $d$ paarweise verschiedenen
gültigen Zeugen. Die Anzahl akzeptierender Pfade ist
$\sum_d c_d f(x)^{\underline d}=P(f(x))$; da $P$ fest ist, bleibt
alles polynomiell.
\end{proof}

In dieser Lesart erhalten die Hori-Operationen eine Zählsemantik:
Addition ist disjunkte Vereinigung von Zählkonstruktionen,
Multiplikation koppelt Konstruktionen (die
Linearisierungskoeffizienten $\binom pj\binom qj j!$ zählen die
Überlappungsmuster), $\Delta$ zählt die durch einen zusätzlichen
Zeugen ermöglichten Strukturen
($(n{+}1)^{\underline d}-n^{\underline d}=d\,n^{\underline{d-1}}$),
und $\sigma(P)=\max(d+c_d)$ misst die maximale Kombination aus
Zeugenanzahl und Variantenanzahl eines Bausteins
(\emph{Zählbreite}; geometrisch die kleinste Treppe über dem
Koeffizientenprofil). Jedes $\mathcal F_r$ ist damit eine endliche
Familie von genau $(r{+}1)!$ \#P-Abschlussoperatoren, und das
Horizontkalkül beziffert optimal das Wachstum der Zählbreite unter
$+$, $\times$, $\Delta$ und $\circ$. Abgrenzung: Die
Binomial-Basis-Vermutung von Ikenmeyer--Pak (deren univariate
Fassung P$\,\ne\,$NP implizieren würde) bleibt unberührt; die
Hori-Operatoren liegen auf der bekannten, relativierenden Seite.
Neu ist ausschließlich die quantitative Hierarchie.

\section{Neun Zählfolgen}

Die folgenden Folgen entstammen der Horimetrik und lieferten in der
OEIS keinen Treffer (zuletzt geprüft am 14.\,Juli 2026; ein
Negativbefund ist kein Neuheitsbeweis). Die Einreichung läuft;
erweiterte Datentafeln (Sprungstellen bis $g=20$, Springer-,
Nicht-Springer- und Fixpunktfolge bis $n=40$, Interchange bis
$H_{12}$, 78 Ein-Ziffern-Terme) liegen als b-Files im
Repositorium:

\begin{center}
\begin{tabular}{ll}
\toprule
Folge & Anfangswerte\\
\midrule
$\beta(r,r)$, maximale Strukturhöhe von Produkten &
$1,6,22,87,407,2473,19527,\dots$\\
Defekt $D(n)$ (Ein-Schritt-Schalenspringer) &
$1,8,60,360,2520,21168,211680$\\
beidseitige Fixpunkte $F(n)=n\cdot n!-D(n)$ &
$4,17,88,540,3960,32760$\\
Sprungstellen $m_g$ des Tiefengesetzes &
$3,15,84,533,3883,31998,294875$\\
Interchange-Maxima $\max\Omega$ auf $H_n$ &
$0,1,6,43,351,2873,26443,270291$\\
Interchange-Nullzähler $\#\{\Omega=0\}$ &
$2,5,13,23,61,111,274,564$\\
Nicht-Springer $n!-D(n)$ &
$5,16,60,360,2520,19152$\\
Ein-Ziffern-Fakultäten $\{n: n!\ \text{einziffrig in } H_{\mu(n!)}\}$ &
$1,2,5,7,11,23,29,39,59,119$\\
Kompositionshöhe $\Gamma_\circ(r,r)$ &
$1,4,23,6541,477149751,\dots$\\
\bottomrule
\end{tabular}
\end{center}

Dabei ist $\Omega(h)=\val(LZh)-\val(ZLh)$ der Defekt zwischen den
beiden Wachstumsreihenfolgen; die Positivität $\Omega\ge0$ ist
bewiesen (Multiplikationslemma mit scharfer Regime-Grenze; zusätzlich
exhaustiv mitbestätigt bis $H_{12}$, also über $6{,}2\cdot10^9$
Elemente), die
Gleichheitsfälle sind charakterisiert und vollständig in
Ziffernbedingungen entfaltet (pro Ebene zwei Ungleichungen bei
wachsendem Exzess), und der Nullzähler gehorcht einer daraus
abgeleiteten Zählformel (exakt für alle gemessenen Terme) mit
bewiesener oberer Schranke $(n{+}1)^{n+1}/n!$ — Gleichgültigkeit
gegenüber der Wachstumsreihenfolge ist asymptotisch extrem selten. Für
$\beta(r,r)$ ist die Asymptotik bewiesen:
$\ln(\beta(r,r)/r!)=2\sqrt r+O(\log r)$ (Reduktion auf einen
Maximalterm plus elementare Schranken). $D(n)$ zählt nach dem Schalensprung-Lemma
genau die Elemente, die eine einzige führende Null in die nächste
Fakultätsschale hebt; die zugehörige Ziffernformel läuft über die
Eigen-Darstellung von $n!$, deren Ein-Ziffern-Fälle vollständig
charakterisiert sind ($n=(i{+}1)!/d-1$, $1\le d\le i$). Alle neun Folgen werden bei der OEIS
eingereicht; die A-Nummern werden nach Vergabe in diesem Text
nachgetragen, sodass Text und Datenbank wechselseitig
referenzieren.

\section{Einordnung und Ehrlichkeit}

Die Bausteine sind bekannte Mathematik: Mischbasensysteme und
Faktoradics, Lehmer-Codes, fallende Fakultäten und endliche
Differenzen, direkte Limites, Produkthalbringe. Die nächsten
bekannten Nachbarn der Nullen-Semantik sind abstrakte
Numerationssysteme; dort ändern führende Nullen den Wert über die
Radix-Ordnung, ohne dass eine Halbringtheorie darüber errichtet
wird. Was wir in dieser Form nicht gefunden haben, ist die
Konstellation: endliche Fakultätshorizonte mit laufender Auswertung,
Wert-/Struktur-Dualität mit zwei Projektoren zur selben Eins, und
ein Starrheitssatz, der die Mischung der Identitäten verbietet.

Die philosophische Motivation --- Unendlichkeit als Treppe endlicher
Horizonte --- ist Motivation, kein Argument. Die belastbare Fassung
lautet: Unendliche Strukturen lassen sich als Systeme endlicher
Horizonte auffassen, und \emph{verschiedene Konsistenzbegriffe
zwischen den Horizonten erzeugen verschiedene globale Räume}
(Satz~\ref{satz:limits}).

Am präzisesten fasst der Grundraum selbst die Lage. Ergänzt man die
Ziffern um eine führende Null, $e=(0,a_1,\dots,a_n)$, so gilt
$0\le e_j<j$: $H_n$ ist als Menge der Raum der Inversionsfolgen der
Länge $n{+}1$, also $H_n\cong S_{n+1}$. Der Strukturhorizont wird
zur Permutationsstatistik
$\sigma(P_a)=n-\min_{e_j>0}((j{-}1)-e_j)$ mit exakter Verteilung
$|\{\sigma=r\}|=r\cdot r!$; ob sie einer klassischen Statistik
entspricht, ist offen. Koordinatenweise wird jeder Horizont zum
distributiven Gitter (die \emph{middle order} auf Permutationen,
arXiv:2405.08943), in dem $\mathcal F_r$ das Hauptideal unter $M_r$
ist. Die Bewohner sind also bekannte Permutationscodes; neu sind die
Operationen darauf --- die umgedrehte Wertung, die Horizontübergänge,
die Dualität, die Seitentreue und das Kapazitätskalkül. Das ist die
günstigste Lage für eine junge Theorie: bekannter Grundraum, neue
Sätze.

Mehrere frühere Vermutungen dieses Projekts wurden durch eigene
Rechnungen widerlegt (eine Schrankenvermutung fiel zweimal an zu
kleinen Testbereichen; eine Formelvermutung an ihrem vierten
Datenpunkt). Wir werten das als Qualitätsmerkmal des Verfahrens,
nicht als Schwäche des Gegenstands.

\section{Offene Fragen}

(1)~Die \emph{oberen} Schwellen des Tiefengesetzes exakt (bisher
nur „spätestens ab“-Schranken; die untere Seite ist vollständig
bewiesen). (2)~Feinstruktur der $\beta$-Asymptotik (vermutet:
$-\tfrac32\ln r$ im Logarithmus-Term). (3)~Allgemeine
Ziffernstatistik von $n!$ in der eigenen Schale (die Ein-Ziffern-Fälle
sind vollständig charakterisiert); daraus die Verteilungsasymptotik
von $D(n)$. (4)~Klassifikation der \emph{gekoppelten} Arithmetiken
innerhalb einer Seite (jenseits der Produktkonstruktionen).
(5)~Eine untere Schranke für den Interchange-Nullzähler.
(6)~$\Gamma_T$ für weitere Operatoren (Komposition, Shift) und die
Frage, ob die Kapazitätshöhe $\max_d(d+c_d)$ anderswo bereits
untersucht wurde. (7)~Formalisierung des Seitentreue-Satzes (die
fallenden Fakultäten liegen in Lean~4/Mathlib bereits vor).

\bigskip
\noindent\textbf{Reproduzierbarkeit.} Alle Enumerationen und
Verifikationen dieses Textes sind als Python-Referenz\-implementierung
mit Selbsttests verfügbar; jedes Resultat trägt in den Arbeitsdateien
den Status bewiesen/vermutet/widerlegt mit Datum. Quellcode,
b-Files, das ausführliche Buch (deutsch und englisch, mit allen
Beweisen) und dieser Text liegen unter
\texttt{github.com/christianbuerckert/horimetrik}; eine archivierte
Fassung liegt auf Zenodo (doi:10.5281/zenodo.21400755).

\begin{thebibliography}{9}
\bibitem{gkp} R.~L. Graham, D.~E. Knuth, O.~Patashnik:
\emph{Concrete Mathematics}, 2.~Aufl., Addison--Wesley, 1994.
\bibitem{rigo} M.~Rigo, M.~Stipulanti: \emph{Revisiting regular
sequences in light of rational base numeration systems},
arXiv:2103.16966.
\bibitem{kiselman} G.~Kudryavtseva, V.~Mazorchuk: \emph{On
Kiselman's semigroup}, arXiv:math/0511374; O.~Ganyushkin,
V.~Mazorchuk: \emph{On Kiselman quotients of 0-Hecke monoids},
arXiv:1006.0316.
\bibitem{ybe} D.~Simon: \emph{Mixed radix numeration bases:
Horner's rule, Yang-Baxter equation and Furstenberg's conjecture},
arXiv:2405.19798.
\bibitem{ikenmeyerpak} C.~Ikenmeyer, I.~Pak: \emph{What is in \#P
and what is not?}, FOCS 2022, arXiv:2204.13149.
\bibitem{middleorder} M.~Bouvel, L.~Ferrari, B.~E.~Tenner:
\emph{Between weak and Bruhat: the middle order on permutations},
arXiv:2405.08943.
\end{thebibliography}

\end{document}
